\subsection{Толковый словарь}
\begin{description}
  \item[Автопилот (\texttt{ap})] программный интерфейс управления полётом: команды взлёта, посадки, полёта к точке, обновления курса.
  \item[События (\texttt{Ev.*})] сигналы от платформы: \texttt{TAKEOFF\_COMPLETE} (взлёт завершён), \texttt{POINT\_REACHED} (точка достигнута), \texttt{COPTER\_LANDED} (посадка), \texttt{SHOCK} (удар/столкновение), \texttt{LOW\_VOLTAGE2} (низкое напряжение), \texttt{ENGINES\_DISARM} (выключение двигателей), \texttt{MCE\_*} (команды режима).
  \item[Обработчик событий (\texttt{callback(event)})] функция, которую платформа вызывает при наступлении системного события.
  \item[Таймер (\texttt{Timer})] механизм неблокирующих вызовов: \texttt{Timer.new(period, fn)} — периодический, \texttt{Timer.callLater(delay, fn)} — отложенный.
  \item[Локальная система координат] плоскость X–Y и высота Z относительно текущей базы; единицы — метры.
  \item[Курс (\texttt{yaw})] угловое направление, выражаемое в радианах; конвертация градусов: \texttt{rad = deg * math.pi / 180}.
  \item[Траектория] набор точек или правило формирования положения во времени (например, окружность).
  \item[Конечный автомат] разбиение логики на состояния с переходами по событиям; состояние хранится в переменной.
  \item[RC-каналы (\texttt{Sensors.rc})] чтение тумблеров/органов управления для запуска миссии или изменения режима.
  \item[Индикаторы (Ledbar)] световая индикация состояния миссии, аварий и этапов.
\end{description}
